\section{运行结果}

\subsection{3量子比特目标态\texorpdfstring{\(|101\rangle\)}{|101>}示例}

选取\(n=3\)、目标态\(|101\rangle\)、迭代次数为2。此时搜索空间大小为\(N=8\)，初始每个基态概率幅为
\[
\frac{1}{\sqrt{8}}\approx 0.353553,
\]
每个基态概率为\(1/8=0.125\)。

\begin{table}[H]
\centering
\caption{\(n=3\)、目标态\(|101\rangle\)时目标态概率幅和概率变化}
\begin{tabular}{llll}
\toprule
步骤 & 操作 & 目标态概率幅 & 目标态概率 \\
\midrule
S0 & Hadamard均匀叠加 & \(+0.353553\) & \(0.125000\) \\
S1 & 第1次Oracle相位翻转 & \(-0.353553\) & \(0.125000\) \\
S2 & 第1次Diffusion均值倒置 & \(+0.883883\) & \(0.781250\) \\
S3 & 第2次Oracle相位翻转 & \(-0.883883\) & \(0.781250\) \\
S4 & 第2次Diffusion均值倒置 & \(+0.972272\) & \(0.945312\) \\
\bottomrule
\end{tabular}
\end{table}

从表中可以看出：
\begin{itemize}
  \item Oracle步骤只改变目标态概率幅符号，不改变目标态概率；
  \item Diffusion步骤把目标态概率幅大幅推高；
  \item 两次完整迭代后，目标态概率达到约\(94.53\%\)；
  \item 如果继续迭代，目标概率会越过最佳点并下降，因此Grover算法需要选择合适的迭代次数。
\end{itemize}

\begin{figure}[H]
  \centering
  \begin{subfigure}[t]{0.48\linewidth}
    \centering
    \IfFileExists{../outputs/figures/00_initial_iter0_amplitude.png}{
      \includegraphics[width=\linewidth]{../outputs/figures/00_initial_iter0_amplitude.png}
    }{
      \fbox{\parbox[c][4.0cm][c]{0.92\linewidth}{\centering 图像文件尚未生成}}
    }
    \caption{初始均匀叠加态的概率幅分布}
  \end{subfigure}
  \hfill
  \begin{subfigure}[t]{0.48\linewidth}
    \centering
    \IfFileExists{../outputs/figures/01_oracle_iter1_amplitude.png}{
      \includegraphics[width=\linewidth]{../outputs/figures/01_oracle_iter1_amplitude.png}
    }{
      \fbox{\parbox[c][4.0cm][c]{0.92\linewidth}{\centering 图像文件尚未生成}}
    }
    \caption{第1次Oracle后目标态概率幅被翻转}
  \end{subfigure}

  \vspace{0.8em}

  \begin{subfigure}[t]{0.48\linewidth}
    \centering
    \IfFileExists{../outputs/figures/02_diffusion_iter1_amplitude.png}{
      \includegraphics[width=\linewidth]{../outputs/figures/02_diffusion_iter1_amplitude.png}
    }{
      \fbox{\parbox[c][4.0cm][c]{0.92\linewidth}{\centering 图像文件尚未生成}}
    }
    \caption{第1次Diffusion后目标态概率幅被放大}
  \end{subfigure}
  \hfill
  \begin{subfigure}[t]{0.48\linewidth}
    \centering
    \IfFileExists{../outputs/figures/03_oracle_iter2_amplitude.png}{
      \includegraphics[width=\linewidth]{../outputs/figures/03_oracle_iter2_amplitude.png}
    }{
      \fbox{\parbox[c][4.0cm][c]{0.92\linewidth}{\centering 图像文件尚未生成}}
    }
    \caption{第2次Oracle后目标态概率幅再次翻转}
  \end{subfigure}

  \vspace{0.8em}

  \begin{subfigure}[t]{0.48\linewidth}
    \centering
    \IfFileExists{../outputs/figures/04_diffusion_iter2_amplitude.png}{
      \includegraphics[width=\linewidth]{../outputs/figures/04_diffusion_iter2_amplitude.png}
    }{
      \fbox{\parbox[c][4.0cm][c]{0.92\linewidth}{\centering 图像文件尚未生成}}
    }
    \caption{第2次Diffusion后目标态概率幅继续增大}
  \end{subfigure}
  \hfill
  \begin{subfigure}[t]{0.48\linewidth}
    \centering
    \IfFileExists{../outputs/figures/target_probability_curve.png}{
      \includegraphics[width=\linewidth]{../outputs/figures/target_probability_curve.png}
    }{
      \fbox{\parbox[c][4.0cm][c]{0.92\linewidth}{\centering 图像文件尚未生成}}
    }
    \caption{目标态概率随完整Grover迭代次数变化}
  \end{subfigure}
  \caption{\(n=3\)、目标态\(|101\rangle\)时两次Grover迭代中的概率幅变化}
\end{figure}

\subsection{与理论结果对照}

对\(N=8\)，有
\[
\theta=\arcsin\left(\frac{1}{\sqrt{8}}\right).
\]
代入理论公式：
\[
P_1=\sin^2(3\theta)=0.78125,
\]
\[
P_2=\sin^2(5\theta)=0.9453125.
\]
这与程序中statevector计算得到的结果一致，说明电路构造和图表数据提取是正确的。

\subsection{2量子比特情况}

当\(n=2\)时，\(N=4\)。单目标Grover搜索在一次完整迭代后即可把目标态概率提升到1。
当\(n=3\)时，目标态概率在多次迭代中先上升后回落，程序通过目标态概率曲线记录这一过迭代现象。
